\documentclass{article}
\usepackage{algorithm, algorithmic, amsfonts, amsmath}

\begin{document}
\begin{algorithm}[!h]
\renewcommand{\algorithmicrequire}{\textbf{Input:}} % show \REQUIRE as Input
\renewcommand{\algorithmicensure}{\textbf{Output:}} % show \ENSURE as Output
	\caption{GPL (graph PU learning with label propagation loss)} % algorithm title
	\begin{algorithmic}[1] % [1] numbers every line
	    \REQUIRE
	    Labeled positive training nodes $P$ and unlabeled training nodes $U$;
	    class prior estimate $\pi_{p}$; original adjacency matrix $A$;
	    outer iterations $E$; inner steps $K$
	    \ENSURE
	    Trained GNN classifier $f_{w}$
	    \STATE Initialize a training GNN classifier $f_{w}$ with $A$
	    \STATE $\hat{A}\leftarrow A$
	    \FOR{training iteration $t=1,2,\dots,E$}
	    \WHILE{$k<K$ or not converged}
	    \STATE $\hat{A}=\arg\min_{\hat{A}}\mathcal{L}_{LPL}(\hat{A})^{(k)}$
	    \hfill \emph{inner loop: reduce heterophilic edges} % heterophilic edges are downweighted; A_hat is a learnable positive mask matrix
	    \ENDWHILE
	    \STATE $\hat{\pi}_{p}=E_{CPE}(\hat{A})$
	    \hfill \emph{estimate $\pi_{p}$ with the optimized $\hat{A}$} % prior estimation is itself disturbed by heterophilic structure, so it comes after heterophily is reduced
	    \STATE Minimize $\mathcal{L}_{GNN}(\hat{\pi}_{p},w,\hat{A})$
	    \hfill \emph{outer loop: train $f_{w}$ with $\hat{A}$ and $\hat{\pi}_{p}$}
	    \ENDFOR
		\RETURN $f_{w}$
	\end{algorithmic}
\end{algorithm}
% Wu, Yao, Han, Yao, Liu, ICML 2024 (arXiv 2405.19919), "Unraveling the Impact of
% Heterophilic Structures on Graph Positive-Unlabeled Learning".
% Compiled from Algorithm 1 and equations (7), (8) and (9) of the original.
% Problem: mainstream PU methods assume instances are generated independently, which fails on graphs, where nodes are
% coupled by edges. The original identifies two ways heterophilic edges, those joining nodes of different classes,
% damage graph PU: (i) they entangle the latent features of positive and negative nodes, and (ii) they make the latent
% labels of unlabeled nodes harder to infer, which in turn degrades class-prior estimation.
% The remedy is a bilevel optimization: the inner level learns an edge-weight matrix A_hat through a label-propagation
% loss (LPL) that suppresses heterophilic edges, and the outer level estimates the prior and trains the GNN on a fixed
% A_hat. The order matters: reduce heterophily, then estimate the prior, then train the classifier, because prior
% estimation is itself affected by heterophilic structure and the error is amplified if the order is reversed.
% LPL depends only on the observed positive nodes, which is what makes it usable under PU (ordinary label propagation
% needs both classes). Convergence follows the standard bilevel results, using the inner solution unrolled J steps.
\end{document}
